%% ============================================================
%% 天津理工大学本科毕业设计说明书（毕业论文）—— 主文件
%% 编译方式：XeLaTeX（本地：latexmk main.tex；Overleaf：Compiler 选 XeLaTeX）
%%
%% 本主文件只负责装配各部分；实际内容写在 chapters/ 下的各文件中。
%% 本模版只覆盖：标题/摘要/关键词、目录、正文、参考文献、附录、致谢。
%% 封皮、首页、原创性声明、选题审批表、任务书、开题报告、底图
%% 使用教务处统一模板，单独制作后合并装订。
%%
%% Copyright (c) 2026 seatres (https://github.com/seatres)
%% 依据 LaTeX Project Public License 1.3c（或更新版本）发布；许可全文见 LICENSE。
%% ============================================================
\documentclass{tjutthesis}

%% 图片搜索路径：插图放进 figures/ 后，\includegraphics 可直接写文件名
\graphicspath{{figures/}}

%% 元数据（届别等可调项）
%% ============================================================
%% 元数据：在此修改届别等可调项（会自动应用到页眉等处）
%% ============================================================

%% 届别：用于页眉“天津理工大学××××届本科毕业设计说明书（毕业论文）”
\gradyear{2026}


\begin{document}

%% ============================================================
%% 前置部分：标题、摘要、关键词、目录（无页眉、无页码）
%% ============================================================
\frontmatter
%% ============================================================
%% 中英文标题 / 摘要 / 关键词（规范二7）
%% 须在 \frontmatter 之后、\tableofcontents 之前 \input
%% ============================================================

%% ---------- 中文标题 / 摘要 / 关键词 ----------
%% 标题不超过 36 个汉字；两人以上课题才需要副标题，
%% 单人课题去掉方括号部分即可：\begin{cnabstract}{标题}
\begin{cnabstract}[实验研究与动力学模型]{考虑电机特性的凸轮分度机构动力学研究}
本文针对凸轮分度机构进行了动力学理论和实验研究，针对机构特点建立了
机构多自由度动力学模型。在此基础上分析了系统的固有特性与动态响应，
并通过实验验证了模型的正确性。研究结果表明……（中文摘要一般在
300--800 字以内，不加注释和评论，应具有独立性和自含性。）
\end{cnabstract}
\cnkeywords{分度凸轮\quad 动力学\quad 模型}

%% ---------- 英文标题 / 摘要 / 关键词 ----------
%% 与中文摘要不另起新页，紧接着写即可；每个关键词首字母大写
\begin{enabstract}[Experimental Research and Dynamic Model]{On the Dynamics of Indexing Cam Mechanisms with Motor Characteristics}
Analytical and experimental research on the dynamics of indexing cam
mechanisms is presented in this thesis. A multi-degree-of-freedom dynamic
model is established according to the characteristics of the mechanism\ldots
\end{enabstract}
\enkeywords{Indexing Cam\quad Dynamics\quad Model}
      % 中英文标题、摘要、关键词
\cleardoublepage
\tableofcontents

%% ============================================================
%% 正文（页码从 1 连续编排；页眉自此开始直到致谢结束）
%% 新增章节：在 chapters/ 下建一个 .tex，再在此处加一行 \input 即可。
%% ============================================================
\mainmatter
%% ============================================================
%% 第一章 绪论
%% ============================================================
\chapter{绪论}
%% 章标题自动排成“第一章 绪论”（三号黑体居中，每章自动另起一页）

本章介绍选题的理由、课题主要解决的问题，采用的手段、方法，简述研究
成果及其意义。正文汉字为小四号宋体，英文、数字为 Times New Roman，
行距 1.25 倍，每个段落首行自动缩进两个汉字。

正文中引用参考文献应在引用处加方括号注明文献号码，作为上角标，如某某人
对此作了研究\cite{ref1}；也可作为文字段落的组成部分，如数学模型见文献
\cite{ref2}。

\section{研究背景与意义}
%% 节标题自动排成“1.1 研究背景与意义”（小三号黑体居中）

这里是节下的正文内容。

\subsection{国内研究现状}
%% 小节标题自动排成“1.1.1 国内研究现状”（四号黑体左起）
%% 标题层级最多到三级（1.1.1），不要再细分

这里是小节下的正文内容。

%% 小节内的小标题：黑体小四单列一行、缩进两个汉字，序号用 1、2、3……自行书写
\minisection{1. 机构动力学等效集中参数模型}

将机构连续参数（质量和刚度连续分布）离散化，用有限多个自由度表示。

%% 段落标题：宋体加黑小四，置于段首；段内序号依次用⑴⑵⑶、A B C、a b c
（1）\paratitle{建模假设}\quad 建立动力学模型过程中作如下假设：
A.~凸轮机构各构件无制造安装误差；B.~凸轮为刚体；C.~考虑凸轮输入轴扭转
弹性，但输入转速为常数。

\subsection{国外研究现状}

这里是小节下的正文内容。

\section{本文主要研究内容}

这里是节下的正文内容。
  % 第一章 绪论
%% ============================================================
%% 第二章 凸轮分度机构的实验研究（图 / 表 / 公式示例）
%% ============================================================
\chapter{凸轮分度机构的实验研究}

\section{实验系统的组成}

图应尽可能置于某页的开始或结尾，并且在图之前的文字段落中有
“如图~x.x”的字样，如图~\ref{fig:demo} 所示。图中的标注一律采用中文，
图题（五号楷体）紧接图的下一行居中打印，图号按章顺序编号。

\begin{figure}[htbp]
  \centering
  %% 把图片放进 figures/ 目录后，直接写文件名即可（已设 \graphicspath）：
  %%   \includegraphics[width=0.6\textwidth]{你的图.png}
  \includegraphics[width=0.5\textwidth]{example-image}
  \caption{实验系统组成示意图}
  \label{fig:demo}
\end{figure}

\section{实验数据与分析}

表应尽可能置于某页的开始或结尾，并且在表之前的文字段落中有
“如表~x.x”的字样，如表~\ref{tab:demo} 所示。表标题（五号楷体）置于
表的上方居中，表内文字为五号；表内不用“同上”“同左”等词，空白代表
无此项内容。

\begin{table}[htbp]
  \centering
  \caption{实验测试结果}
  \label{tab:demo}
  \begin{tabular}{cccc}
    \toprule
    序号 & 转速/(r/min) & 扭矩/(N·m) & 温升/℃ \\
    \midrule
    1 & 1000 & 12.5 & 8.2 \\
    2 & 1500 & 13.1 & 10.6 \\
    3 & 2000 & 13.8 & 13.4 \\
    \bottomrule
  \end{tabular}
\end{table}

数学、物理和化学式在正文中另起一行打印，式的序号按章顺序编排并标注在
该行最右边，如式~(\ref{eq:demo})：
\begin{equation}
  m\ddot{x} + c\dot{x} + kx = F(t)
  \label{eq:demo}
\end{equation}
式中：$m$ 为等效质量；$c$ 为阻尼系数；$k$ 为刚度；$F(t)$ 为激励力。
较长的式另行居中横排；如必须转行，只能在 $+$、$-$、$\times$、$\div$、
$<$、$>$ 处转行，上下式尽可能在等号“$=$”处对齐。
    % 第二章 凸轮分度机构的实验研究
%% ============================================================
%% 第三章 结论
%% ============================================================
\chapter{结论}

对本人所做工作进行归纳和综合，得到设计或研究的结论。与已有结果进行
比较，指明所得结论的新进展，并对该课题尚应进一步改进与研究的方向提出
建议。文字要简单、明确。
    % 第三章 结论

%% ============================================================
%% 文末部分
%% ============================================================
%% ============================================================
%% 参考文献（自动另起一页；“前引后列”：每条都须在正文中被 \cite 引用，
%% 并按引用先后顺序排列；著录格式按 GB/T 7714-2015）
%% ============================================================
\begin{thebibliography}{99}

%% 期刊文献 [J]：主要责任者. 题名[J]. 期刊名，年，卷(期):页码.
\bibitem{ref1} 张三，李四，王五，等. 凸轮分度机构动力学建模与仿真研究[J]. 机械工程学报，2024，60(3):45-52.

%% 学位论文 [D]：主要责任者. 题名[D]. 大学所在城市:大学名称，出版年.
\bibitem{ref2} 赵六. 高速凸轮机构动态特性分析[D]. 天津:天津理工大学，2023.

%% 普通图书 [M]：主要责任者. 书名[M]. 版本项. 出版地:出版者，出版年:引文页码.
\bibitem{ref3} 孙七. 机械原理[M]. 3版. 北京:机械工业出版社，2021:120-135.

%% 电子资源 [EB/OL]：必须列出资料所在的具体页面、引用日期与访问路径
\bibitem{ref4} 周八. 工业凸轮机构设计手册[EB/OL]. (2023-05-12)[2026-03-01]. https://www.example.com/cam-handbook.

%% 外文期刊（作者超过三位列前三位，后加 et al.）
\bibitem{ref5} Smith J, Brown K, Davis M, et al. Dynamic analysis of indexing cam mechanisms[J]. Mechanism and Machine Theory, 2023, 180: 105-118.

\end{thebibliography}
       % 参考文献
%% ============================================================
%% 附录（不宜放正文但有参考价值的内容：较复杂的公式推演、程序等；
%% 没有附录可在 main.tex 中删去对应的 \input 行）
%% ============================================================
\begin{appendices}

\chapter{主要程序清单}
%% 自动编为“附录1 主要程序清单”（三号黑体居中）
%% 附录中文为 5 号宋体，英文数字 5 号 Times New Roman，行距固定值 18 磅

这里放置较复杂的公式推演、计算机程序等内容。

\end{appendices}
         % 附录（无附录可删去本行）
%% ============================================================
%% 致谢（自动另起一页，页眉到此结束）
%% ============================================================
\begin{acknowledgements}
以简短的文字对指导教师和协助完成设计（论文）的人员表示谢意。
感谢×××老师在毕业设计期间给予的悉心指导……
\end{acknowledgements}
  % 致谢

\end{document}
